% 图契约：电容充电时能量从极板边缘流入极板间电场区域；几何、方向和公式与正文同一约定。
\documentclass[tikz,border=5pt]{standalone}
\usepackage{fontspec,xeCJK,amsmath,amssymb}
\IfFontExistsTF{Times New Roman}{\setmainfont{Times New Roman}}{\setmainfont{TeX Gyre Termes}}
\IfFontExistsTF{Heiti SC}{\setCJKmainfont[AutoFakeBold=1.5]{Heiti SC}}{\setCJKmainfont[AutoFakeBold=1.5]{FandolHei-Regular}}
\usetikzlibrary{arrows.meta,calc,patterns}
\definecolor{ink}{HTML}{243745}
\definecolor{blue}{HTML}{275DAB}
\definecolor{green}{HTML}{177C63}
\definecolor{red}{HTML}{BD4935}
\definecolor{orange}{HTML}{AA6C16}
\definecolor{purple}{HTML}{715393}
\begin{document}
\begin{tikzpicture}[x=1cm,y=1cm,font=\fontsize{9.5}{12}\selectfont,
v/.style={-{Latex[length=2mm,width=1.3mm]},line width=.85pt},
axis/.style={-{Latex[length=1.7mm,width=1.1mm]},ink!65,line width=.55pt},
aux/.style={ink!45,dashed,line width=.55pt},
lab/.style={fill=white,inner sep=1.8pt},
card/.style={draw=ink!25,fill=blue!3,rounded corners=2pt,text width=6.1cm,minimum height=.85cm,inner sep=6pt,align=center}]
\path[use as bounding box](-8,-3.6)rectangle(8,3.65);
\node[font=\bfseries\fontsize{10.5}{13}\selectfont,text width=15.1cm,align=center]at(0.00000,3.22000){M7-H06　电容充电时能量从极板边缘流入极板间电场区域};
\begin{scope}[shift={(-4,0)}]
\fill[blue!6](-2,-1.5)rectangle(2,1.5);
\draw[ink,line width=1.6pt](-2.00000,-1.50000)--(-2.00000,1.50000);
\node[]at(-2.00000,1.85000){$+$};
\draw[ink,line width=1.6pt](2.00000,-1.50000)--(2.00000,1.50000);
\node[]at(2.00000,1.85000){$-$};
\draw[v,blue](-1.70000,-1.00000)--(1.70000,-1.00000);
\draw[v,blue](-1.70000,-0.33000)--(1.70000,-0.33000);
\draw[v,blue](-1.70000,0.33000)--(1.70000,0.33000);
\draw[v,blue](-1.70000,1.00000)--(1.70000,1.00000);
\node[lab,text=blue]at(0.00000,0.60000){$\mathbf E$};
\draw[{Latex[length=1.3mm]}-{Latex[length=1.3mm]},ink!65,line width=.5pt](-2.00000,-2.00000)--(2.00000,-2.00000);
\node[lab]at(0.00000,-2.25000){$d$};
\node[]at(0.00000,2.35000){板面积 $A$，介质 $\varepsilon$};
\draw[v,red](-1.00000,2.00000)--(-1.00000,1.30000);
\draw[v,red](-1.00000,-2.00000)--(-1.00000,-1.30000);
\draw[v,red](1.00000,2.00000)--(1.00000,1.30000);
\draw[v,red](1.00000,-2.00000)--(1.00000,-1.30000);
\draw[green](0,1.8)circle(0.13);
\fill[green](0,1.8)circle(.035);
\node[anchor=west,text=green]at(0.22000,1.80000){$\mathbf H$};
\node[text=red]at(0.00000,-2.60000){能流从边缘进入板间};
\end{scope}
\node[card]at(4.00000,1.25000){$\partial_tu_E>0$};
\node[card]at(4.00000,-1.05000){$\mathbf S\cdot\mathbf n<0$：充电储能};
\end{tikzpicture}\end{document}
